\documentclass[english]{article}

\usepackage{amssymb,amsmath,amsfonts,amstext,amsthm}
\usepackage{bbm}
\newcommand{\bx}{\textbf{x}}
\newcommand{\bv}{\textbf{v}}
\newcommand{\bB}{\textbf{B}}
\newcommand{\bE}{\textbf{E}}
\newcommand{\bJ}{\textbf{J}}
\newcommand{\du}{\,\mathrm{d}}

\newcommand{\nr}{n_{\mathrm{ref}}}
\newcommand{\qr}{q_{\mathrm{ref}}}
\newcommand{\Tr}{T_{\mathrm{ref}}}
\newcommand{\mr}{m_{\mathrm{ref}}}
\newcommand{\Br}{B_{\mathrm{ref}}}
\newcommand{\betar}{\beta_{\mathrm{ref}}}
\newcommand{\Lr}{L_{\mathrm{ref}}}
\newcommand{\vr}{v_{\mathrm{ref}}}
\newcommand{\omr}{\omega_{\mathrm{ref}}}
\newcommand{\rhor}{\rho_{\mathrm{ref}}}
\newcommand{\epsr}{\varepsilon_{\mathrm{ref}}}

\newcommand{\phiht}{\hat{\phi}}
\newcommand{\psigaht}{\hat{\psiga}}
\newcommand{\Apht}{\hat A_{\parallel}}
\newcommand{\Tsht}{\hat{T_{s}}}
\newcommand{\fht}{\hat{f}}
\newcommand{\qht}{\hat{q}}
\newcommand{\mht}{\hat{m}}
\newcommand{\Eht}{\hat{\bE}}
\newcommand{\Bht}{\hat{\bE}}
\newcommand{\rhoht}{\hat{\rho}}
\newcommand{\jht}{\hat{\bJ}}
\newcommand{\fsht}{\hat{f_s}}
\newcommand{\tht}{\hat{t}}
\newcommand{\xht}{\hat{\bx}}
\newcommand{\vht}{\hat{\bv}}
\newcommand{\muht}{\hat{\mu}}
\newcommand{\vpht}{\hat{\vp}}

\title{Code units in GEMPIC}

\author{Katharina Kormann}

\begin{document}

\maketitle

Vlasov equation for the distribution function of s particle species $s$ reads in SI units
\begin{equation}\label{Vlasov}
\partial_t f_s + \bv \cdot \nabla_{\bx} f_s + \frac{q_s}{m_s} \left( \bE + \bv \times \bB \right) \cdot \nabla_{\bv} f_s = 0
\end{equation}
and Maxwell's equations
\begin{subequations}
\begin{align}
\label{ampere}
\nabla \times \bE &= - \frac{\partial \bB}{\partial t}, \\
\label{electric_gauss}
\nabla \times \bB &= \mu_0 \left( \bJ + \varepsilon_0 \frac{\partial \bE}{\partial t} \right), \\
\label{faraday}
\nabla \cdot \bE &= \frac{\rho}{\varepsilon_0},\\
\label{magnetic_gauss}
\nabla \cdot \bB &= 0.
\end{align}
\end{subequations}
The source terms of the Maxwell's equation are given by the velocity moments of the distribution function: the charge density
\begin{align*}
\rho= \sum_s q \int f_s \du\bv
\end{align*}
and the current density
\begin{align*}
\bJ =\sum_s q\int \bv f_s \du\bv.
\end{align*}
The total energy of the system is given by the Hamiltonian in SI units, 
\begin{align}\label{hamiltonian}
\mathcal{H}=\sum_s\frac{m_s}{2}\int |\bv|^2 f_s \du\bx \du\bv + \frac{\varepsilon_0}{2} \int \|\bE\|^2 \du\bx + \frac{1}{2\mu_0} \int \|\bB\|^2 \du\bx.
\end{align}

In the code, we use a normalization that prescribes the following reference values:
\begin{itemize}
	\item $\nr$, $\mr$, $\qr$ 
	\item two out of  $\omr$, $\Lr$, $\vr$. 
\end{itemize}
There is the following relation between the reference quanties:
\begin{itemize}
\item The reference velocity $\vr = \sqrt{\frac{\Tr k_B}{\mr}}$,
\item The  reference length or frequency $\Lr = \frac{\vr}{\omr}, \omr = \frac{\vr}{\Lr}$,
\item The reference epsilon $\epsr = \frac{\qr^2\nr}{\omr^2\mr}$,
\item The reference beta $\betar = \mu_0 \varepsilon_0 \vr^2$.
\end{itemize}


There are two common choices:
\begin{itemize}
	\item $\omega_{c,i}$ (gyrofrequency), $d_i$, $v_{A,i}$ (Alfven velocity), $\betar = \beta_i$ (plasma beta)
	\item $\omega_{p,i}$ (plasma frequency), $d_i$ ($d_i$ ion inertial length), $c$, $\epsr = \varepsilon_0$
\end{itemize}


The dimensionless Vlasov--Maxwell system takes the following form: 
\begin{align*}
\partial_{\tht} \fsht + \vht \cdot \nabla_{\xht} \fsht + \frac{\qht}{\mht} \left( \Eht + \vht \times \Bht \right) \cdot \nabla_{\vht} \fsht &= 0,
\end{align*}

\begin{align*}
\frac{\partial \Bht}{\partial \tht}&=-\nabla_{\xht} \times \Eht,   &\nabla_{\xht}  \cdot \Eht &= \frac{\epsr}{\varepsilon_0}\rhoht,
\\ 
\frac{\partial \Eht}{\partial \tht} &=\frac{1}{\betar }\nabla_{\xht} \times \Bht -\frac{\epsr}{\varepsilon_0 } \jht, 
&\nabla_{\xht} \cdot \Bht  &= 0 
\end{align*}

with the source terms
\begin{align*}
\rhoht= \sum_s  \hat q_{s} \int  \fsht \du\vht,  \quad  \jht=\sum_s \hat q_{s}  \int \vht \fsht \du\vht.
\end{align*}


The dimensionless Hamiltonian takes the following form:
\begin{align*}
\hat{\mathcal{H}}= \sum_s \frac{\mht_s}{2}\int  |\vht|^2  \fsht \du\xht \du\vht + \frac{\varepsilon_0 }{\epsr}\frac{1}{2} \int  \|\Eht \|^2 \du\xht  + \frac{1}{\mu \epsr \vr^2 } \frac{1}{2} \int \|\Bht \|^2 \du\xht.
\end{align*}


\end{document}